% ---------------------------------------------------------------------
% 2.12 Reducción de canales y repeticiones
% ---------------------------------------------------------------------
\section[REDUCCIÓN DE CANALES Y REPETICIONES]{Reducción de canales y repeticiones}

Como se señaló en la introducción, los sistemas P300 basados en EEG suelen alcanzar una alta
exactitud a costa de utilizar muchos electrodos y de repetir varias veces la estimulación de cada
fila y columna. Ambos factores afectan directamente la practicidad del sistema: un mayor número de
canales incrementa el costo del equipo y el tiempo necesario para colocar los electrodos, además de
dificultar su uso fuera del laboratorio, mientras que un mayor número de repeticiones reduce la
velocidad de comunicación y aumenta la fatiga del usuario \cite{krusienski2008,contreras2021}.

\subsection{Selección de canales}

La selección de canales busca identificar el subconjunto mínimo de electrodos que conserve la
información necesaria para detectar la P300. Krusienski et al. evaluaron distintos conjuntos de
electrodos y encontraron que un conjunto de ocho canales, Fz, Cz, Pz, P3, P4, PO7, PO8 y Oz, que
combina posiciones de la línea media con posiciones parieto-occipitales, ofrecía un desempeño
superior al de los tres canales de la línea media utilizados tradicionalmente \cite{krusienski2008}.
Por otra parte, Cecotti y Gräser aprovecharon los pesos aprendidos por la capa espacial de su CNN
para ordenar los electrodos según su contribución a la clasificación y eliminar los menos
relevantes \cite{cecotti2011}, estrategia que resulta natural en arquitecturas como EEGNet, cuya
convolución \textit{depthwise} aprende explícitamente un peso por electrodo \cite{lawhern2018}. En
el extremo opuesto, Contreras trabajó con solo cuatro electrodos de la línea media y obtuvo valores
de AUC aceptables en la evaluación intrasujeto \cite{contreras2021}.

\subsection{Reducción de repeticiones}

De acuerdo con la Ecuación~\eqref{eq:promediado}, cada repetición adicional mejora la relación
señal-ruido de la P300 en un factor proporcional a $\sqrt{N}$, pero también añade alrededor de
2.1~s al tiempo de cada selección con los parámetros de la base de datos utilizada. Una primera
estrategia para reducir las repeticiones consiste en mejorar la clasificación de ensayo único, es
decir, la capacidad del modelo de detectar la P300 en una sola época sin promediarla con otras,
objetivo que persiguen las arquitecturas convolucionales descritas en la Sección~\ref{sec:cnn}
\cite{cecotti2011,lawhern2018}. Una segunda estrategia es la parada dinámica, en la que el sistema
deja de presentar repeticiones en cuanto la confianza acumulada sobre algún carácter supera un
umbral, en lugar de usar siempre un número fijo; Speier et al. mostraron que combinar este tipo de
clasificación dinámica con información del lenguaje mejora tanto la exactitud como la tasa de
transferencia del P300 Speller \cite{speier2012}.
